\par\noindent\textbf{(1)} 测量位置后选相应一侧，把该侧容积V/2等温可逆扩张至V，单分子状态方程p=k\_BT/V'，输出功W\_ext=∫\_\{V/2\}\textasciicircum{}V(k\_BT/V')dV'=k\_BT ln2。气体内能在理想等温过程不变，热库供热恰等于该功。
\par\noindent\textbf{(2)} 左右结果等概率，测量获得的Shannon信息I=−Σp\_i ln p\_i=ln2（单位nat）。信息反馈热力学在最优协议下满足⟨W\_ext⟩≤k\_BTI，因此本模型达到k\_BT ln2；若测量有错误，互信息小于ln2，平均可提取功界随之降低。
\par\noindent\textbf{(3)} 把对两种结果随机占据的记忆重置为一个固定逻辑0，记忆熵至少下降k\_Bln2。与温度T热库接触时Landauer原理要求向环境至少排热Q\_erase≥k\_BTln2，只有准静态、可逆复位可达到等号。
\par\noindent\textbf{(4)} 完整循环需将记忆、控制器和活塞恢复初态，提取功不超过复位最低费用，整体净功不为正且总熵不减少。真实误差、有限速操作与测量装置耗散提高总成本；单独忽略记忆复位才会产生第二类永动机表象。
